\documentclass[a4paper,11pt]{article}
\usepackage{fontspec}
\setmainfont{CMU Serif}
\setmonofont{CMU Typewriter Text}
\newfontfamily\cyrillicfonttt{CMU Typewriter Text}
\usepackage{polyglossia}
\setmainlanguage{russian}
\usepackage[margin=20mm]{geometry}
\usepackage{graphicx,array}
\title{Проектирование корпоративного практикума по конкурентному программированию на Go}
\author{Адаменко С. С.}
\date{2026}
\begin{document}
\maketitle
\begin{abstract}
Предлагается структура практикума для разработчиков, изучающих синхронизацию и взаимодействие горутин. Учебный результат проверяется самостоятельным решением новой задачи. Текст представляет проект фрагмента ВКР; результаты эксперимента не заявляются.
\end{abstract}
\section{Постановка задачи}
Конкурентность связана с организацией независимо выполняющихся действий, параллелизм - с их одновременным выполнением. Обучающийся должен обнаруживать несогласованный доступ к общим данным и обосновывать исправление. Модель памяти Go описывает условия видимости изменений между горутинами. Общие изменяемые данные требуют синхронизации с помощью каналов или примитивов пакетов \texttt{sync} и \texttt{sync/atomic} \cite{mem}.
\section{Структура модуля}
Этапы показаны на рисунке~\ref{fig:module}. После проверки возможно повторное обращение к практике.
\begin{figure}[h]
\centering
\setlength{\unitlength}{1mm}
\begin{picture}(160,18)
\put(0,2){\framebox(48,14){Диагностика}}
\put(55,2){\framebox(48,14){Практика}}
\put(110,2){\framebox(48,14){Перенос навыка}}
\end{picture}
\caption{Схема учебного модуля}\label{fig:module}
\end{figure}
\begin{table}[h]
\centering
\caption{Задания и наблюдаемые результаты}\label{tab:tasks}
\begin{tabular}{|p{32mm}|p{63mm}|p{50mm}|}\hline
Этап & Действие & Результат \\\hline
Диагностика & Найти несогласованный доступ & Объяснение ошибки \\\hline
Практика & Исправить синхронизацию & Код и тесты \\\hline
Перенос & Решить новую задачу & Обоснование решения \\\hline
\end{tabular}
\end{table}
\newpage
\section{Проверка и обратная связь}
Детектор гонок запускается командой \texttt{go test -race}. Он выявляет гонки в выполненных участках программы; отсутствие сообщений не доказывает корректность всех выполнений \cite{race}. Автоматическая проверка дополняется разбором жизненного цикла горутин и условий завершения.

Предлагается оценивать корректность синхронизации, качество тестов и объяснение решения. Каждый критерий получает от нуля до двух баллов:
\begin{equation}
S=C+T+E,\quad C,T,E\in\{0,1,2\},\quad 0\leq S\leq6.
\label{eq:score}
\end{equation}
Формула~\ref{eq:score} задаёт авторскую шкалу, а не валидированный инструмент. Таблица~\ref{tab:tasks} связывает действия с результатами.
\subsection{Корпоративный контекст}
Задания воспроизводят рабочие ситуации, но используют обезличенный учебный код. Внешним сервисам не передаются закрытые исходные тексты компании. Подсказки постепенно сокращаются; итоговое задание выполняется самостоятельно.
\section{Вывод}
Практикум связывает техническую проверку с оценкой понимания. Следующий этап - пилотное проведение курса и проверка переноса навыка в профессиональные задачи.
\begin{thebibliography}{9}
\bibitem{mem} The Go Authors. The Go Memory Model. URL: \texttt{https://go.dev/ref/mem} (31.10.2026).
\bibitem{race} The Go Authors. Data Race Detector. URL: \texttt{https://go.dev/doc/articles/race\_detector} (31.10.2026).
\end{thebibliography}
\end{document}